\documentclass[14pt,a4paper]{extarticle}
\usepackage[T2A]{fontenc}
\usepackage[utf8]{inputenc}
\usepackage[english,russian]{babel}
\usepackage{lmodern}
\usepackage{geometry}
\usepackage{array}
\usepackage{tabularx}
\usepackage{longtable}
\usepackage{ltablex}
\usepackage{booktabs}
\usepackage{hyperref}
\usepackage{indentfirst}
\usepackage{microtype}
\usepackage{graphicx}
\usepackage{setspace}
\usepackage{listings}
\usepackage{graphicx}
\usepackage{caption}
\usepackage{float}
\usepackage{tocloft}
\usepackage{titlesec}
\usepackage{newtxtext,newtxmath}

\titleformat{\section}        % форматируем \section
  {\normalfont\bfseries}      % стиль: обычный + жирный
  {\thesection.}              % номер секции с точкой
  {2mm}                       % расстояние между номером и текстом
  {}                           % перед заголовком можно добавить что-то
\titlespacing*{\section}{1.5cm}{*3}{*1}  % {отступ слева}{отступ сверху}{отступ снизу}

\titleformat{\subsection}        % форматируем \section
  {\normalfont\bfseries}      % стиль: обычный + жирный
  {\thesubsection.}              % номер секции с точкой
  {2mm}                       % расстояние между номером и текстом
  {}                           % перед заголовком можно добавить что-то
\titlespacing*{\subsection}{1.5cm}{*3}{*1}  % {отступ слева}{отступ сверху}{отступ снизу}

\captionsetup{justification=centering}
\renewcommand{\ttdefault}{cmtt}
\keepXColumns
\setlength{\parindent}{3em}
\newgeometry{left=2.5cm,right=2.5cm,top=3cm,bottom=2.5cm}

\renewcommand{\contentsname}{\hspace*{4em}{\smallСодержание}}
\setlength{\cftbeforesecskip}{0.5em}
\renewcommand{\cftsecleader}{\cftdotfill{\cftdotsep}}
\renewcommand{\cftsecaftersnum}{.\hspace{2mm}}

\title{Отчёт по "Управление ИТ-проектами"}
\author{Выполнил студент группы М-РИТ-25-1, Станиславчук Сергей}

\begin{document}

\setstretch{1.5}
\noindent
\begin{minipage}[t]{0.01\textwidth}
    \includegraphics{~/gitrepos/LSTU_MAG/5_2/LSTU.jpg}
\end{minipage}%
\begin{minipage}[c]{1.1\linewidth}
    \centering
    {\small
    МИНИСТЕРСТВО НАУКИ И ВЫСШЕГО ОБРАЗОВАНИЯ РФ \\[1em]
    ФЕДЕРАЛЬНОЕ ГОСУДАРСТВЕННОЕ БЮДЖЕТНОЕ \\
    ОБРАЗОВАТЕЛЬНОЕ УЧРЕЖДЕНИЕ ВЫСШЕГО ОБРАЗОВАНИЯ \\
    «ЛИПЕЦКИЙ ГОСУДАРСТВЕННЫЙ ТЕХНИЧЕСКИЙ УНИВЕРСИТЕТ»
    }
\end{minipage}%

\vspace{3em}

\noindent
\hspace{1em}%
\begin{minipage}[t]{0.8\textwidth}
\textbf{Институт} \hspace{2em} компьютерных наук \\
\textbf{Кафедра} \hspace{2em} автоматизированных систем управления
\end{minipage}


\vspace{6em}

\begin{center}
\textbf{Практическая работа №12} \\
по дисциплине \\ <<Управление ИТ-проектами>>
\end{center}

\vspace{6em}

\noindent
\begin{tabular*}{\textwidth}{@{\extracolsep{\fill}} l c r }
    Студент М-РИТ-25-1 & \makebox[4cm]{\hrulefill} & Станиславчук С.М. \\[2em]
Руководитель & \\
К.т.н., доцент & \makebox[4cm]{\hrulefill} & Ведищев В.В. \\
\end{tabular*}

\vfill
\begin{center}
    Липецк, 2026
\end{center}

\thispagestyle{empty}

\newpage
\tableofcontents
\newpage

\section{Управление сборками и релизами программного продукта}

\subsection{Понятие «сборки»}

Сборка (build) представляет собой процесс преобразования исходного кода системы в исполняемый или готовый к развертыванию программный продукт. В рамках разрабатываемой системы корпоративного обмена сообщениями сборка включает компиляцию серверной и клиентской частей, установку зависимостей, выполнение автоматических тестов и формирование артефактов (исполняемых файлов, контейнеров или пакетов развертывания).

Сборки могут выполняться автоматически при внесении изменений в систему контроля версий. Это позволяет оперативно выявлять ошибки интеграции и поддерживать работоспособное состояние проекта на каждом этапе разработки.

\subsection{Понятие «релиза»}

Релиз представляет собой зафиксированную и протестированную версию программного продукта, предназначенную для использования в рабочей среде. В отличие от сборки, релиз проходит полный цикл проверки, включая функциональное, нагрузочное и, при необходимости, безопасностное тестирование.

В рамках проекта релиз формируется на основе стабильной версии системы, находящейся в основной ветке разработки. Каждому релизу присваивается уникальный номер в соответствии с принятой схемой версионирования, а также формируется набор сопровождающей документации, включая описание изменений и инструкции по развертыванию.

\subsection{Распределение ролей в команде при выпуске релиза}

Процесс формирования релиза требует координации нескольких ролей. Бэкенд-разработчик отвечает за подготовку серверной части, контроль интеграции изменений и корректность формирования версии продукта. Он также принимает решение о готовности системы к выпуску.

Фронтенд-разработчик обеспечивает готовность клиентского приложения, проверяет корректность пользовательского интерфейса и совместимость с серверной частью.

Тестировщик выполняет проверку релизной версии, включая функциональное и регрессионное тестирование. По результатам тестирования принимается решение о выпуске или доработке системы.

Специалист по информационной безопасности оценивает корректность реализации механизмов защиты, проверяет наличие уязвимостей и соответствие системы требованиям безопасности.

Системный администратор отвечает за развертывание релиза в целевой среде, настройку серверов и обеспечение стабильной работы системы после выпуска.

\subsection{Инструментальные средства формирования релизов}

Для управления сборками и релизами в проекте используется система контроля версий Git, обеспечивающая хранение исходного кода, ветвление и отслеживание изменений.

Основным инструментом сборки является CMake, применяемый для кроссплатформенной генерации файлов сборки. CMake позволяет описать структуру проекта, зависимости между модулями, параметры компиляции и линковки. На его основе генерируются make-файлы, используемые утилитой Make для непосредственного выполнения сборки.

Процесс сборки включает конфигурацию проекта (этап генерации), компиляцию исходного кода, линковку исполняемых файлов и формирование итоговых артефактов. Использование Make обеспечивает эффективную инкрементальную сборку, при которой пересобираются только изменённые компоненты.

Автоматизация сборки и тестирования может быть реализована через CI/CD-инструменты (например, GitLab CI), где сборка проекта выполняется с использованием CMake и Make. Это позволяет автоматически проверять корректность изменений, запускать тесты и формировать релизные версии при каждом обновлении репозитория.

Для подготовки релизов могут формироваться архивы исполняемых файлов или пакеты, готовые к развёртыванию в корпоративной инфраструктуре. Дополнительно используются средства логирования и мониторинга для контроля работы системы после внедрения.

Таким образом, применение CMake и Make обеспечивает воспроизводимость сборки, контроль зависимостей и упрощает процесс формирования стабильных релизов программного продукта.

\end{document}
